% 《大学物理学学习指导》第5版，自测题9（光的干涉）选择题9.1.1—9.1.11
% 原书印刷202—204页，PDF物理210—212页；原图依文字方向重述。
\subsection{自测题9：干涉选择题}

\sourcetest{9.1.1}{202}{210}
\problemtext 单色光真空波长为$\lambda$，在折射率$n$的媒质中
从A传播到B，两点光相位差为$3\pi$，求该段光程。
\workings
光程与相位差的关系为
$\Delta\phi=2\pi\mathcal L/\lambda$，故
\[
\boxed{\mathcal L=\frac{\Delta\phi}{2\pi}\lambda=\frac32\lambda}.
\]
原书选\textbf{A}。
\crosscheck 问的是光程$n\,AB$，不是几何长度$AB=3\lambda/(2n)$。

\sourcetest{9.1.2}{202}{210}
\problemtext 图9.1.1：透射薄膜折射率$n_2$、厚$e$，上下环境
$n_1<n_2>n_3$；入射光在介质$n_1$中的波长为$\lambda_1$。
两束反射光的光程差是多少？
\workings
上表面反射从疏到密，出现一次$\pi$相变；
下表面从密到疏无半波相变。
真空波长$\lambda_0=n_1\lambda_1$，反射两路的
有向有效光程差可记为
\[
\boxed{\Delta=2n_2e-\frac{n_1\lambda_1}{2}},
\]
对应原书选\textbf{C}。
\crosscheck 相对半波损失符号取决于两束相减顺序；
反射强弱只依模$2\pi$的相位差，不能遗漏$n_1\lambda_1=\lambda_0$。

\sourcetest{9.1.3}{202}{210}
\problemtext 图9.1.2：杨氏双缝光源S从两缝等距离位置向下移至S'，
观察中央明纹的位置变化及条纹间距。
\workings
S下移使下缝S$_2$相对S$_1$的入射光程缩短。
要恢复零光程差，屏上观察点需位于上方，使
S$_2$到该点的传播路径更长。
故中央明纹\textbf{上移}。
两缝间距、波长、屏距均不变，
$\Delta x=\lambda D/d$仍不变，选\textbf{B}。
\crosscheck 这里只移动上游光源而没有改变双缝结构，
造成的是初相差的平移而非条纹密度变化。

\sourcetest{9.1.4}{202}{210}
\problemtext 白光照射双缝，分别用纯红和纯蓝滤光片遮住两缝，
判断屏上的干涉现象。
\workings
来自不同颜色的两缝光波频率不同；
长期平均的交叉干涉项
$\langle E_R E_B\rangle$随时间平均为零，
因此没有两缝间稳定的红蓝相干条纹，
\[
\boxed{\text{不形成稳定双缝干涉条纹}},
\]
选\textbf{D}。
\crosscheck 可能见到各单缝的衍射照度分布，
不等于红蓝光之间会形成两套相干双缝干涉条纹。

\sourcetest{9.1.5}{202--203}{210--211}
\problemtext 图9.1.3：双缝原先在P点形成亮纹。
遮住S$_2$直接光后，将镜M放在其附近使反射光到P，
判断P点亮暗。
\workings
新光路与原S$_2$直达光路的几何长度一般不同，
而反射镜还引入与材料相关的反射相位。
若未进一步给镜的位置及光程差，不能确定
与S$_1$直达光在P处的相位是偶数还是奇数$\pi$。
故只能断言P点\textbf{不一定仍亮或变暗}，
选\textbf{C：无法确定}。
\crosscheck 即使理想金属镜近似产生固定反射相位，
光路的额外几何长度仍不可忽略。

\sourcetest{9.1.6}{203}{211}
\problemtext 空气劈尖的左端为棱边，上玻璃板绕该棱边
按原图约定逆时针微转，使劈角增大，判断反射等厚纹变化。
\workings
空气劈尖厚度$t(x)\simeq\alpha x$，
同级厚度$t_m$与角$\alpha$无关，
$x_m=t_m/\alpha$；邻纹距
\[
\Delta x=\frac{\lambda}{2\alpha}.
\]
$\alpha$增加使$\Delta x$变小，
同级纹向左侧棱边移动。
选\textbf{A（变密并向棱边移动）}。
\crosscheck 转角方向以源题的上板与左侧棱边图示为准；
若画图镜像，方向口述必须随图调整。

\sourcetest{9.1.7}{203}{211}
\problemtext 牛顿环透镜垂直向上缓慢移动而远离下面的平板玻璃，
问观察到的反射干涉环怎样变化。
\workings
空气膜厚度$t(r)=t_0+r^2/(2R)$，
固定级数暗环$2t(r)=m\lambda$，故
\[
r_m^2=R(m\lambda-2t_0).
\]
上移使$t_0$增大，固定级数的$r_m$减小，
环纹\textbf{向中心收缩}，选\textbf{B}。
\crosscheck 当$m\lambda<2t_0$时该级暗环消失，
不能继续取负的半径平方。

\sourcetest{9.1.8}{203}{211}
\problemtext 空气中独立放置折射率$n$透明薄膜，
波长$\lambda$的光垂直入射；反射光增强的最小正膜厚。
\workings
空气$\to$膜有一次半波相变，
膜$\to$空气反射无半波相变。
反射增强要求$2nt=(m+1/2)\lambda$，
取$m=0$，
\[
\boxed{t_{\min}=\frac{\lambda}{4n}},
\]
选\textbf{B}。
\crosscheck 此结果只针对\emph{空气--薄膜--空气}，
不能直接套用于膜下还有高折射率基底的减反膜问题。

\sourcetest{9.1.9}{203}{211}
\problemtext 玻璃基底$n_3=1.60$上镀MgF$_2$膜$n_2=1.38$，
空气$n_1=1.00$中500nm光垂直入射，求反射尽量弱时最小正膜厚。
\workings
由于$n_1<n_2<n_3$，两主要反射均有$\pi$相变，
相对反射相位损失抵消。
相消条件$2n_2t=\lambda/2$，所以
\[
\boxed{t_{\min}=\frac{500}{4(1.38)}\,\mathrm{nm}
\simeq90.6\,\mathrm{nm}},
\]
选\textbf{E}。
\crosscheck 只消除相对相位并不足以保证两个反射振幅相等；
因此题面准确说法是“尽可能少反射”。

\sourcetest{9.1.10}{203}{211}
\problemtext 用劈尖干涉探查工件表面的局部凹凸，
原图9.1.5显示直条纹在缺陷区发生一个条纹间距的横向弯曲。
判断缺陷方向与高度/深度。
\workings
同一干涉级次要求空气膜厚度不变，
且相邻反射等厚纹的膜厚差为$\lambda/2$。
原图局部弯曲相当于一个条纹间距，
所以缺陷表面相对正常面有
\[
\boxed{|h|=\frac{\lambda}{2}}.
\]
根据原图劈尖由左到右增厚的斜率，
局部干涉纹向薄侧弯曲表示空气膜在缺陷处更厚，
因此工件表面应有\textbf{凹陷}，
选\textbf{C（凹陷深度$\lambda/2$）}。
\crosscheck 图中“条纹向哪侧弯”决定凹/凸符号。
如果原图扫描倒置或劈尖厚度方向被重画反转，
须重新逐线判定；相位级差$\lambda/2$不受此影响。

\sourcetest{9.1.11}{203--204}{211--212}
\problemtext 迈克耳孙干涉仪一条臂插入折射率$n$、
厚度$d$的薄片，求这一臂光程相对不插片时的改变。
\workings
光线经过该薄片向镜面行进一次并返回一次，
每次较空气中的同厚光程增加$(n-1)d$，所以
\[
\boxed{\Delta\mathcal L=2(n-1)d},
\]
选\textbf{A}。
\crosscheck 迈克耳孙仪是往返光路，不能漏掉因子2。
